\documentclass[12pt]{article}
\usepackage{ceai}

\begin{document}

\title{CE 488/5588 Homework 05 - Concrete Strength Classification Using Machine Learning\\[2in]
\textbf{Name: \rule{2in}{0.15mm}}}
\author{}
\date{}
\maketitle

\newpage
\doublespacing

\section{Introduction}

This project extends the concepts introduced in Tutorials 04-A and 04-B on supervised classification using the UCI Concrete Compressive Strength dataset.

You will implement and evaluate machine learning classifiers that distinguish between normal-strength concrete (NSC) and high-strength concrete (HSC). The project emphasizes both computational analysis and engineering interpretation.

The tasks involve three stages:

\begin{enumerate}
\item Defining the classification problem and establishing a baseline model
\item Training a more expressive classifier and evaluating its capability
\item Performing external validation using your own designed concrete mixtures
\end{enumerate}

All implementations should be carried out using Python in a Google Colab environment.

\section{Dataset and Classification Definition}

The dataset used in this project is the UCI Concrete Compressive Strength dataset.

Each sample includes mixture composition variables such as cement, slag, fly ash, water, and age, along with the measured compressive strength.

You should define a binary classification label separating NSC and HSC.

\subsection{Classification Threshold}

In the tutorials, a nominal threshold of

\[
f'_c(28) = 40 \text{ MPa}
\]

was used.

For this project, you may choose your own nominal threshold value for defining the classes.

The classification rule should be written as

\[
y =
\begin{cases}
1 & f'_c(28) \le T \\
0 & f'_c(28) > T
\end{cases}
\]

where

\begin{itemize}
\item $T$ is your chosen threshold
\end{itemize}

Please briefly justify your chosen value of $T$.

\section{Baseline Model}

You should construct a baseline classifier using either

\begin{itemize}
\item Logistic regression, or
\item Perceptron
\end{itemize}

\subsection{Model Description}

Please briefly describe the model used, including

\begin{itemize}
\item the decision function
\item the activation function
\item the classification rule
\end{itemize}

\subsection{Training and Testing}

Using a train-test split procedure, please report

\begin{itemize}
\item training accuracy
\item testing accuracy
\end{itemize}

\subsection{Confusion Matrix}

Please compute and report the confusion matrix for the testing dataset.

From this matrix, you should compute the following performance measures:

\begin{itemize}
\item Accuracy
\item Precision
\item Recall
\item F1 score
\end{itemize}

\section{Advanced Model}

You should train a more expressive classifier using one of the following approaches:

\begin{itemize}
\item Multi-Layer Perceptron (MLP) with at least two hidden layers
\item Support Vector Machine (SVM) with a nonlinear kernel
\end{itemize}

\subsection{Model Design}

Please specify the model architecture or configuration.

For an MLP, please describe

\begin{itemize}
\item number of hidden layers
\item number of neurons per layer
\item activation functions
\end{itemize}

For an SVM, please specify

\begin{itemize}
\item kernel type (e.g., polynomial or RBF)
\item kernel parameters
\item regularization parameter $C$
\end{itemize}

You should briefly explain your design choices.

\subsection{Training and Testing Performance}

Please compute and report the following metrics for the advanced model:

\begin{itemize}
\item training accuracy
\item testing accuracy
\item confusion matrix
\item precision
\item recall
\item F1 score
\end{itemize}

\section{Classifier Capability Analysis}

To evaluate classifier capability beyond simple accuracy, you should analyze performance across different classification thresholds.

\subsection{ROC Curve}

Please generate the ROC curve for both

\begin{itemize}
\item the baseline model
\item the advanced model
\end{itemize}

Both curves should be plotted on the same figure.

Please also compute the area under the curve (AUC).

\subsection{Precision–Recall Curve}

Please generate the precision–recall curve for both models and overlay them on the same figure.

\subsection{Model Capability Comparison}

You should compare the performance of the two models and discuss:

\begin{itemize}
\item Which model appears more capable?
\item Does the nonlinear model improve classification performance?
\item Are there trade-offs between precision and recall?
\end{itemize}

\section{External Validation Using Designed Mixtures}

In this section, you will evaluate your trained model using concrete mixtures designed by yourself.

\subsection{Designing Concrete Mixtures}

Please design two concrete mixtures:

\begin{itemize}
\item one intended to produce normal-strength concrete (NSC)
\item one intended to produce high-strength concrete (HSC)
\end{itemize}

For each mixture, please report:

\begin{itemize}
\item cement content
\item water content
\item supplementary cementitious materials
\item aggregate quantities (if included)
\item curing age
\end{itemize}

You should briefly explain the reasoning behind your mixture design.

\subsection{Model Prediction}

Using your trained classifier, please compute the predicted class for each designed mixture.

Please report:

\begin{itemize}
\item predicted probability
\item predicted class
\end{itemize}

\subsection{Engineering Interpretation}

Please discuss whether the prediction aligns with your expected strength category.

If discrepancies occur, please comment on possible reasons, such as

\begin{itemize}
\item limitations of the training dataset
\item mixture variables not captured in the dataset
\item uncertainty in the strength estimation model
\end{itemize}

\section{Project Deliverables}

Please submit:

\begin{itemize}
\item a short report (PDF)
\item the Colab notebook used for analysis
\end{itemize}

Your report should include:

\begin{itemize}
\item description of the classification threshold
\item description of the baseline model
\item description of the advanced model
\item training and testing performance metrics
\item ROC and PR curve comparisons
\item mixture design and external validation results
\item discussion and conclusions
\end{itemize}
\end{document}
